\bookchapter{Глава 28: Покупаем с умом — финальное руководство по выбору смартфона}

\chapterintro{Покупка телефона — дело, к которому многие подходят основательно: неделю смотрят обзоры, месяц сравнивают цены, а в итоге всё равно спонтанно тратят пять-шесть тысяч юаней на флагман, который им по сути не нужен. В этой главе не будет таблиц с параметрами — их и так можно найти на сайте производителя. Я хочу дать вам «мышление при выборе», чтобы в дальнейшем, какие бы новинки ни появлялись на рынке, вы могли сами решить, подходит вам этот телефон или нет.}

\booksection{Ноль. Метод трёх шагов: потребности бюджет модель}

Многие в интернете спрашивают: «Какой телефон купить за три тысячи?» На самом деле на этот вопрос ответить нельзя — потому что неизвестно, для чего вы покупаете телефон. Правильный порядок размышлений при покупке должен быть таким:

\begin{tcbinfo}{Три топора выбора}
\begin{enumerate}
  \item \textbf{Сначала определите, что для вас важнее всего}: камера? игры? автономность? экран? плавность системы? или достаточно просто звонков и WeChat? Расставьте приоритеты.
  \item \textbf{Затем установите верхнюю границу бюджета}: не «примерно столько-то», а «максимум, который я готов потратить». Жёсткий лимит не даст смотреть всё более дорогие варианты.
  \item \textbf{И только потом выбирайте модель}: с вашим списком приоритетов и бюджетом подбирайте конкретные устройства.
\end{enumerate}

\end{tcbinfo}

\booksection{Один. На что смотреть в характеристиках — не нужно знать всё, достаточно пяти пунктов}

В таблицах характеристик десятки параметров, но обычному пользователю достаточно обращать внимание на пять измерений:

\begin{center}
\begin{tabularx}{\linewidth}{|X|X|X|}
\hline
\textbf{Параметр} & \textbf{На что смотреть} & \textbf{Итог в одном предложении} \tabularnewline
\hline
Процессор (SoC) & Определяет, сколько лет телефон будет работать плавно. Стандарт 2026 года: Snapdragon 8 Elite Gen5 / Dimensity 9500 и выше = топ; Snapdragon 8s Gen4 / Dimensity 8400 = средне-высокий; Snapdragon 7+ Gen4 / Dimensity 7350 = хватит на три года. Ниже Snapdragon 6-й серии — подходить с осторожностью. & Покупайте новое, а не старое — SoC определяет, сколько лет телефон не будет тормозить \tabularnewline
Экран & OLED лучше LCD по цветам и экономичнее; частота 120 Гц плавнее 60 Гц; разрешение — минимум 1080P. ШИМ-димминг с высокой частотой (от 1920 Гц) меньше утомляет глаза. & OLED + 120 Гц + высокочастотный ШИМ = проходной балл хорошего экрана \tabularnewline
Камера & \textbf{Смотрите не на мегапиксели, а на размер сенсора}. Большой сенсор важнее количества мегапикселей. Сверхширокоугольник полезен; телеобъектив (от 3x оптического зума) — для съёмки удалённых объектов; макрообъектив обычно бесполезен. & Основная камера с большим сенсором + полезный телевик > россыпь «для галочки» \tabularnewline
Батарея и зарядка & 4500 м Ач — проходной минимум, свыше 5000 м Ач — зарядки раз в день достаточно. Мощность зарядки: 60 Вт — быстро, 100 Вт и выше — очень быстро. Беспроводная зарядка удобна, но медленная. & 5000 м Ач + 60 Вт = комбинация без тревоги \tabularnewline
Обещание обновлений & Определяет, как долго телефон будет безопасно использоваться. Apple обычно даёт 5–6 лет системных обновлений; флагманы Samsung — 4 года системы + 5 лет безопасности; Google Pixel — 7 лет; китайские среднебюджетные модели — обычно 2–3 года. & Хотите пользоваться дольше — выбирайте бренды с длительной поддержкой \tabularnewline
\hline
\end{tabularx}
\end{center}

\booksection{Два. Рекомендации по бюджету (на середину 2026 года)}

Рекомендации ниже основаны на рынке материкового Китая 2026 года, цены — официальные стартовые. Реальные цены могут колебаться, ориентируйтесь на ту цену, по которой фактически покупаете.

\booksubsection{До 1000 юаней: звонить, WeChat, и хватит}

\textbf{Кому подходит}: студенты, второй телефон, непритязательные пользователи, для звонков и WeChat.

\textbf{Что можно купить}: Redmi Note, Honor X, realme цифровой серии — младшие комплектации. Процессоры уровня Snapdragon 6-й серии или Dimensity 6000, экран в основном LCD, камера снимет, но фокус медленный, плавность системы средняя.

\textbf{Совет}: в этом сегменте смотрите на две вещи: ёмкость батареи (от 5000 м Ач) и объём памяти (минимум 6 ГБ + 128 ГБ, не берите 4 ГБ — будет тормозить). Redmi в этом сегменте имеет наиболее явное преимущество.

\booksubsection{1000–2000 юаней: лучшее соотношение цена/качество, 3–4 года без лагов}

\textbf{Кому подходит}: студенты, ограниченный бюджет, но с расчётом на долгое использование, без претензий к топовым камере или играм.

\textbf{Типичные модели}: Redmi Turbo, iQOO Z, Honor X старшие комплектации, OnePlus Ace начального уровня.

\textbf{Совет}: выбирайте модели с Snapdragon 7+ Gen3 или Dimensity 7300 и выше, с памятью 8 ГБ + 256 ГБ. Этого хватит на повседневные задачи + эпизодические сессии в Honor of Kings, телефон проработает 3–4 года без тормозов.

\booksubsection{2000–3500 юаней: средний класс, заметный рост производительности и камеры}

\textbf{Кому подходит}: пользователи, которые хотят хорошую камеру, плавные игры и долгий срок службы. Это самый популярный ценовой диапазон.

\textbf{Типичные модели}: Xiaomi цифровой серии (например Xiaomi 17 Standard со скидкой), vivo S старшие, OPPO Reno, Honor цифровой серии, Redmi K.

\textbf{Совет}: главная черта этого сегмента — «универсалы» без явных слабых мест. По сравнению с моделями до 2000 юаней: камера ощутимо лучше (появляется оптическая стабилизация), экран — качественный OLED, быстрая зарядка (часто 80 Вт и выше), качество сборки тоже выше. Берите 12 ГБ + 256 ГБ или 12 ГБ + 512 ГБ — проработает 4–5 лет без проблем.

\booksubsection{3500–5000 юаней: субфлагманы — всем, кроме гиков, более чем достаточно}

\textbf{Кому подходит}: тем, кто ценит качество камеры и топовые материалы корпуса.

\textbf{Типичные модели}: Xiaomi цифровой серии Standard, vivo X, OPPO Find X Standard, прошлогодние iPhone (например iPhone 16).

\textbf{Совет}: в этом диапазоне уже можно взять прошлогодние флагманы. Моя рекомендация:\textbf{лучше прошлогодний флагман, чем нынешний средний сегмент}. У прошлогоднего флагмана материалы, сборка и камера обычно на класс выше, чем у нынешнего среднего. Главное — быть готовым, что это не «самая новая» модель.

\booksubsection{Свыше 5000 юаней: флагманский сегмент, для тех, у кого бюджет не ограничен или есть особые запросы}

\textbf{Кому подходит}: фотоэнтузиасты, заядлые геймеры, ценители топовых экранов и материалов, авторы контента, снимающие видео на телефон.

\textbf{Типичные модели}: iPhone 17 Pro / Pro Max, Samsung S26 Ultra, Huawei Pura 80 Ultra, Xiaomi 17 Ultra, vivo X300 Pro.

\textbf{Совет}: на этом уровне разница между брендами становится всё меньше — выбор в основном по экосистеме. Нравится камера — vivo или Huawei, нравится видеосъёмка — iPhone, нравится экран и S Pen — Samsung, нравится баланс — Xiaomi. Если честно, на этом уровне\textbf{можно выбирать по симпатии — берите тот, что вам по душе}.

\booksection{Три. Выбор по сценарию — какие у вас задачи?}

Ниже — несколько самых распространённых «портретов пользователя», выберите свой:

\begin{center}
\begin{tabularx}{\linewidth}{|X|X|X|}
\hline
\textbf{Сценарий} & \textbf{На что смотреть} & \textbf{Рекомендация} \tabularnewline
\hline
Фотоэнтузиаст & Основная камера с большим сенсором (лучше 1 дюйм), оптическая стабилизация, телеобъектив (от 3x оптического зума), выделенный чип обработки изображений & vivo X, Huawei Pura/Mate, Xiaomi Ultra \tabularnewline
Геймер & Топовый процессор, высокая частота дискретизации касаний, система охлаждения, большая батарея + быстрая зарядка & iQOO цифровой серии, ROG Phone, Red Magic, Redmi K \tabularnewline
Приоритет автономности & Батарея от 5500 м Ач, энергоэффективный процессор, LTPO-экран & Honor Magic, OnePlus, Redmi Note версия с большой батареей \tabularnewline
Использование пожилыми & Упрощённый режим, крупный шрифт, громкий звук, большая батарея, много сервисных центров & Huawei Changxiang, Honor Play, OPPO A, Redmi Note \tabularnewline
Ученик & Хорошее соотношение цена/качество, прочность, большая батарея, наличие родительского контроля & Redmi Note/Turbo, realme GT, iQOO Neo \tabularnewline
Пользователь экосистемы Apple & Если уже есть Mac/iPad/AirPods, iPhone — единственный логичный выбор & iPhone 17 Standard (цена/качество) или Pro (полный опыт) \tabularnewline
Путешествия за границу & Поддержка eSIM, полное покрытие глобальных частот, сервисы Google & iPhone, Samsung S, Pixel, OnePlus \tabularnewline
\hline
\end{tabularx}
\end{center}

\booksection{Четыре. Пять типичных ловушек — проверьте перед покупкой}

\begin{enumerate}
  \item \textbf{Ловушка 1: четыре-пять камер $\neq$ хорошая съёмка}— производители любят набирать количество объективов, среди них часто два — 2-мегапиксельные «для галочки» (макро + глубина), совершенно бесполезные. Для оценки камеры достаточно смотреть на\textbf{основную + сверхширокоугольную + телеобъектив}. Остальное можно игнорировать.
  \item \textbf{Ловушка 2: не ведитесь на 100 мегапикселей}— высокое разрешение не равно хорошему качеству. Размер сенсора, диафрагма и алгоритмы обработки важнее. Многие флагманы с 50-мегапиксельным сенсором снимают гораздо лучше, чем некоторые «сотни мегапикселей» в бюджетных моделях.
  \item \textbf{Ловушка 3: не покупайте операторские версии}— «пополнение счёта + телефон» выглядит выгодно, но такие модели получают обновления медленнее, имеют предустановленные приложения, которые нельзя удалить, и обычно заблокированы под конкретного оператора. Кажется, что экономите, а на деле «привязаны» к нему.
  \item \textbf{Ловушка 4: «изогнутый экран — это премиум» — это маркетинг, а не истина}— изогнутые экраны когда-то были признаком флагмана, но всё больше людей замечают их недостатки: сложно наклеить плёнку, ложные касания по краям, дорогой ремонт при трещине. Тренд 2026 года — возврат к «слегка изогнутым» и даже полностью плоским экранам. Не выбирайте изогнутый экран только потому, что «изогнутый = премиальный».
  \item \textbf{Ловушка 5: не все «AI-телефоны» действительно полезны}— в 2026 году все кричат про «AI-телефон», но многие просто переименовали старый голосовой ассистент. По-настоящему полезные AI-функции: перевод в реальном времени, AI-стирание (убрать человека с фото), AI-ретушь, локальные большие модели, локальные умные саммари. Остальные «фишки» — в основном маркетинг.
\end{enumerate}

\booksection{Пять. Когда покупать выгоднее всего?}

\begin{itemize}
  \item \textbf{618 (конец мая — середина июня)}: самые большие скидки на телефоны в году. Многие производители снижают цену на флагманы на 800–1500 юаней, чтобы выполнить план продаж. Лучшее окно для покупки в первой половине года.
  \item \textbf{11.11 (конец октября — середина ноября)}: большая распродажа второй половины года, масштаб сопоставим с 618. В это время у многих брендов только что вышло новое поколение флагманов, и предыдущее поколение максимально дешевеет.
  \item \textbf{Через месяц после выхода новой модели}: в день выхода — официальная цена, через месяц — каналы (сторонние магазины) обычно дают скидку 200–500 юаней.
  \item \textbf{За две недели до Нового года и до Праздника фонарей}: офлайн-магазины распродают остатки в рамках «праздничной распродажи», можно поторговаться с продавцом.
  \item \textbf{Не рекомендуется}: покупать в день презентации по предзаказу — если только вы не хотите быть первым пользователем, цена в первый день — самая высокая.
\end{itemize}

\booksection{Шесть. Где покупать — как выбрать надёжный канал?}

\begin{itemize}
  \item \textbf{Официальный магазин бренда / Tmall Flagship / JD Self-operated}: самый безопасный вариант, цена не самая низкая, но нет риска восстановленного или поддельного товара, лучший сервис.
  \item \textbf{Pinduoduo «Bai Yi Bu Tie» (субсидия 10 млрд)}: самая низкая цена, но нужно выбирать магазины с пометкой «Bai Yi Bu Tie». Без такой пометки — не берите. При получении запишите видео распаковки — если что-то не так, можно вернуть.
  \item \textbf{Офлайн-магазины}: можно пощупать и попробовать. Учтите: консультанты продвигают модели с высокой маржой (а не те, что подходят вам). Лучше заранее выберите 2–3 целевые модели в интернете, а в магазине сразу говорите: «Я хочу посмотреть только эти».
  \item \textbf{Площадки б/у (Xianyu / Zhuanzhuan)}: дёшево, но рискованно. Если покупаете б/у, помните три правила: требуйте чек или скриншот заказа, требуйте IMEI для проверки,\textbf{пользуйтесь только встроенной системой сделки, никогда не переводите деньги напрямую}.
\end{itemize}

\begin{tcbpractice}{Практика: выберите телефон по этой методике}
\begin{enumerate}
  \item Возьмите лист бумаги, запишите две главные потребности (например: автономность — первое, камера — второе).
  \item Установите верхний предел бюджета (например 3500 юаней).
  \item Откройте JD или Tmall, отфильтруйте по цене, выберите 3–5 моделей-кандидатов.
  \item Зайдите на Bilibili, найдите «модель XX — опыт использования за месяц» (не обзор в первый день, а реальные впечатления после минимум месяца использования), посмотрите пару-тройку видео.
  \item Убедившись в надёжности канала, оформите заказ и запишите видео распаковки при получении.
\end{enumerate}

\end{tcbpractice}

\begin{tcbsummary}{Итоги главы}
\begin{itemize}
  \item Три топора выбора: сначала потребности затем бюджет в конце модель
  \item Смотрите только на пять ключевых параметров: процессор, экран, камера, батарея/зарядка, обещание обновлений
  \item 1500–3500 юаней — самый выгодный и рекомендуемый сегмент
  \item Прошлогодний флагман выгоднее, чем нынешний средний сегмент
  \item Пять ловушек: обман многокамерности, ловушка мегапикселей, операторские модели, мифы об изогнутых экранах, AI-хайп
  \item Распродажи 18 июня (618) и 11 ноября (11.11) — два лучших временных окна в году для покупки телефона
\end{itemize}

\end{tcbsummary}
